
% Adrianic Multi-Slot Complex Memory Candidate

\subsection{Memory Bank}

Let the memory bank contain $M$ independent complex fields:
\begin{align}
\mathcal B
=
\{(\chi_j,\mathbf v_j,u_j,t_j)\}_{j=1}^{M},
\end{align}
where $\chi_j$ is the stored dynamical template, $\mathbf v_j$ is an
interpretable address vector, $u_j$ is usage metadata, and $t_j$ is recency.

\subsection{Transparent Address Vector}

For the current vortex-family implementation:
\begin{align}
\mathbf v
=
(
x_1,y_1,x_2,y_2,
x_c,y_c,d,
\bar x_A,\bar y_A,
\Sigma_{xx},\Sigma_{yy},\Sigma_{xy}
).
\end{align}

The first seven coordinates describe the two low-amplitude core locations,
their center and separation. The remaining five describe amplitude centroid
and covariance.

\subsection{Address Distance}

For per-coordinate scale $s_k$:
\begin{align}
D_j
=
\sqrt{
\frac{1}{K}
\sum_{k=1}^{K}
\left(
\frac{v_k-v_{jk}}{s_k}
\right)^2
}.
\end{align}

Let $D_{(1)}$ and $D_{(2)}$ be the best and second-best matches.
Define the confidence margin
\begin{align}
\Delta_D=D_{(2)}-D_{(1)}.
\end{align}

If $\Delta_D<\Delta_{\min}$, the address controller returns
\[
\boxed{\mathrm{HOLD\ UNRESOLVED}}
\]
rather than forcing a memory selection.

\subsection{Recall}

For selected slot $j$:
\begin{align}
\left.\frac{\partial\psi}{\partial t}\right|_{\rm recall}
=
R_j(t)\mu_R(\chi_j-\psi).
\end{align}

\subsection{Commit and Replacement}

If an experience is accepted for storage and an empty slot exists, allocate it.
If the bank is full, a transparent replacement utility may be
\begin{align}
U_j
=
\frac{n_j}{1+\alpha a_j},
\end{align}
where $n_j$ is use count and $a_j$ is time since last access.
The lowest-utility slot is the replacement candidate.

\subsection{Ambiguity Rule}

If two stored states share the same reliable address information after corruption,
the system must not infer missing information.  For example, opposite chiralities
with identical amplitudes are indistinguishable after total phase erasure unless
a contextual or temporal cue is available.
